\documentclass[10pt,a4paper]{amsart}
\usepackage[margin=19mm]{geometry}
\usepackage{amsmath,amssymb,mathtools}
\usepackage[scr=rsfs,cal=euler]{mathalfa}
\usepackage{xfrac}
\usepackage[unicode,hidelinks,bookmarks=false]{hyperref}
\newcommand{\ie}{i.e.\ }
\newcommand{\cf}{cf.\ }
\newcommand{\resp}{respectively\ }
\newcommand{\Subsection}{\subsection}
\providecommand{\nobreakdash}{\nobreak\hbox{-}}
\setlength{\emergencystretch}{2em}
\multlinegap0pt
\usepackage{amssymb}
\usepackage{xcolor}
\usepackage{tikz}
\usepackage[shortlabels]{enumitem}
\usepackage{float}
\newtheorem{lem}{Lemma}
\newtheorem{thm}{Theorem}
\newtheorem{cor}{Corollary}
\usepackage{cleveref}
\crefname{lem}{Lemma}{Lemmas}
\crefname{thm}{Theorem}{Theorems}
\crefname{cor}{Corollary}{Corollarys}
\newcommand{\N}{\mathbb{N}}
\newcommand{\Soc}{\operatorname{Soc}}
\newcommand{\Stab}{\operatorname{Stab}}
\newcommand{\Z}{\mathbb{Z}}
\newcommand{\id}{\operatorname{id}}
\title{Finiteness conditions on skew braces and solutions of the Yang-Baxter equation}
\author{Rosa Cascella, Silvia Properzi, Arne Van Antwerpen}
\date{Source: arXiv:2603.06177v2 (CC BY 4.0)}
\begin{document}
\maketitle

\noindent\emph{Source and licence.} Finiteness conditions on skew braces and solutions of the Yang-Baxter equation. Version arXiv:2603.06177v2 (CC BY 4.0), distributed by arXiv under the Creative Commons Attribution 4.0 International licence, http://creativecommons.org/licenses/by/4.0/, as recorded in the arXiv licence metadata for this exact version and shown on the abstract page https://arxiv.org/abs/2603.06177v2. Full licence notice: \texttt{licenses/arxiv-2603.06177-CC-BY-4.0.txt}.\par\smallskip

\noindent\textit{Editorial scope.} Counted unit: the article's thm unsatisfacttheorem (source0.tex lines 1161--1165). The article's complete original proof of it is delivered with it (source0.tex lines 1166--1210, one \texttt{proof} environment of 45 lines). No mathematics is written, repaired or re-derived: the statement and the proof are the article's own lines, in the article's own notation, and no retained line is substituted or re-flowed.  Retained uncounted as the source-local context the proof needs: lines 610--618; lines 653--660; lines 897--900; lines 1100--1105; lines 1131--1136; lines 1143--1149.  Nothing was omitted from any retained span, so the package carries no omission record.  No retained line contains a citation, so the wrapper carries no bibliography and no bibliography entry.  No retained line contains a figure environment, an image-inclusion command or drawing code, so no image is delivered and the package declares no asset dependency.  Editorial formatting only, in addition: an AMS \texttt{amsart} preamble that restates the article's own declarations and macros used by the retained lines, namely the environments \texttt{lem}, \texttt{thm}, \texttt{cor} on separate counters, together with the standard packages those lines need.  The retained environments print in this document as Lemma 1, Theorem 1, Corollary 1 in order of appearance, whereas the article prints the same items with the numbering of its own full document.  The complete native source of the article as distributed in the arXiv e-print for this exact version is bundled unchanged as \texttt{sources/195844/source0.tex}.
\begin{lem}
\label{gens}
    Let $B$ be a skew brace and let $T$ be a set of generators (as a skew brace) of $B$.
    Then, denoting $T'=\bigcup_{b\in B}\lambda_b(T)$,
    \[
    U=T' \cup (-T')
    \] 
    generates $(B,+)$ and $(B,\circ)$.
\end{lem}

\begin{thm}\label{fg}
    Let $B$ be a $\lambda_f$-skew brace. Then, the following are equivalent.
    \begin{enumerate}
        \item The skew brace $B$ is finitely generated as a skew brace,
        \item the group $(B,\circ)$ is finitely generated,
        \item the group $(B,+)$ is finitely generated.
    \end{enumerate}
\end{thm}

\begin{lem}
\label{thetafg}
    Let $(B,+,\circ)$ be a skew brace such that the additive group $(B,+)$ is generated by a finite set of $\theta_f$-elements $\{x_1,...,x_t\}$. Then, $|B:\Soc(B)|$ is finite. Moreover, if for all $i$ one has that $|[x_i]_{\theta}| \leq k$, then $|B:\Soc(B)| \leq k^{t}.$
\end{lem}

\begin{cor}
\label{Dietzmann cor}
    Let $B$ be a skew brace and let $x$ be a $\theta_f$-element of finite additive order. Then $x$ has finite multiplicative order.

    Furthermore, if $B$ is a $\theta_f$-skew brace, then the torsion $T_+(B)$ of $(B,+)$ forms a locally finite strong left ideal that is contained in the set of torsion elements of $(B,\circ)$.
\end{cor}

\begin{lem}
\label{BFC -> G/Z has finite exp}
    Let $(G,\cdot)$ be a BFC-group with $k\in\N$ such that
    $|[g]_\cdot|\leq k$ for all $g\in G$.
    Then $G/Z(G)$ is of finite exponent at most $k!$.
\end{lem}

\begin{lem}
\label{order of lambda for theta Bf}
    Let $(B,+,\circ)$ be a $\theta_{Bf}$ skew brace
    with $k\in\N$ such that
    $|[x]_\theta|\leq k$ for all $x\in B$.
    Then $\lambda_b$ has finite order $\leq k!$.
\end{lem}

\begin{thm}\label{unsatisfacttheorem}
    Let $B$ be a $\theta_f$-skew brace.
    Then every $\lambda_b$ is of finite order if and only if $B/\Soc(B)$ is periodic. 
    Moreover, if $B$ is $\theta_{bf}$, then $B/\Soc(B)$ is of finite exponent.
\end{thm}

\begin{proof}
Let $B$ be a $\theta_f$-skew brace and  $x \in B$.
Then, we know that $\Stab_{\lambda}(x) $ is 
of finite index $\leq [x]_\theta=n$
in $(B,\circ)$.
Take $T=\left\lbrace t_1,..., t_n \right\rbrace$ to be a transversal. 
Then, the skew brace $A=\left<t_1,...,t_n,x\right>$ is $\theta_f$ and finitely generated. 
Hence, by \cref{fg} $(A,+)$ is finitely generated. 
More precisely, by \cref{gens}, 
the group $(A,+)$ is generated by 
\[
\bigcup_{i=1}^n\left\lbrack t_i\right\rbrack_{\theta} \cup [x]_{\theta},
\]
which contains at most $n'(n+1)$ elements, 
where $n'=\max\{\max\{|[t_i]_\theta|:i\in \{1,...,n\}\},|[x]_\theta|\}$.
So, by \cref{thetafg},
$\Soc(A)$ is of finite index $m=|A:\Soc(A)|\leq n^{n'(n+1)}$ in $A$.
Thus, $\lambda_{mx}(t_i) = t_i$ and $\lambda_{mx}(x) =x$ and $mx \in Z(A,+)$ for $m\leq n^{n'(n+1)}$.
Since $(B,+)$ is an $FC$-group,
$(B,+)/Z(B,+)$ is periodic.
So  there exists $h\in \Z$ such that $hmx\in Z(B,+)$.
Moreover, $mx\in\Soc(A)$, therefore 
$hmx=(mx)^{\circ h}$ and so 
$\lambda_{hmx}=(\lambda_{mx})^h$
is the identity on $A$.
Therefore $hmx\in \Soc(A)\cap Z(B,+)$.
Since $\lambda_{hmx}$ is of finite order,
there exists a $l$ such that $\lambda_{hmx}^l=\textup{id}_B$.
As $hmx\in \Soc(A)$, we see that $\lambda_{lhmx}=\lambda_{hmx}^l=\id_B$.
As $lhmx \in Z(B,+)$, it follows that $lhmx \in \Soc(B)$, which shows the first half of the result. 

Assume now that $B/\Soc(B)$ is periodic. As $B/\Soc(B)$ is an $\theta_f$-skew brace, \cref{Dietzmann cor} shows that every element of $B/\Soc(B)$ has finite multiplicative order. 
Hence, for every $b \in B$ there exists a $n\in \N$ such that $\lambda_b^n=\id_B$, which shows the result.

In the case of $B$ being $\theta_{Bf}$,
let $N\in \N$ such that $[x]_\theta\leq N$.
Going through the calculations before,
we can see that $n,n'\leq N$, $m\leq N^{N(N+1)}$.
Moreover, since $[x]_+\subseteq [x]_\theta$ for every $x\in B$,
using \cref{BFC -> G/Z has finite exp}, we get that $h\leq N!$.
Finally, by \cref{order of lambda for theta Bf}, $ l\leq N!$.
Therefore we can conclude that $N! N!  N^{N(N+1)} x\in\Soc(B)$
for every $x\in B$.
Hence $B/\Soc(B)$ has finite exponent.
\end{proof}

\end{document}
